\documentclass[10pt,a4paper]{amsart}
\usepackage[margin=19mm]{geometry}
\usepackage{amsmath,amssymb,mathtools}
\usepackage[scr=rsfs,cal=euler]{mathalfa}
\usepackage{xfrac}
\usepackage[unicode,hidelinks,bookmarks=false]{hyperref}
\newcommand{\ie}{i.e.\ }
\newcommand{\cf}{cf.\ }
\newcommand{\resp}{respectively\ }
\newcommand{\Subsection}{\subsection}
\providecommand{\nobreakdash}{\nobreak\hbox{-}}
\setlength{\emergencystretch}{2em}
\multlinegap0pt
\usepackage{bm}
\usepackage{bbm}
\usepackage{dsfont}
\usepackage{xspace}
\usepackage{cancel}
\usepackage{mathtools}
\usepackage{cleveref}
\usepackage{amsthm}
\makeatletter
\@ifundefined{theorem}{\newtheorem{theorem}{Theorem}}{}
\@ifundefined{lemma}{\newtheorem{lemma}[theorem]{Lemma}}{}
\@ifundefined{remark}{\newtheorem{remark}[theorem]{Remark}}{}
\makeatother
\makeatletter
\newcommand{\R}{\mathbb{R}}
\expandafter\ifx\csname abstract\endcsname\relax
\renewenvironment{abstract}{
  \par\smallskip
  \noindent\textsc{Abstract.}\ %
}{
  \par\medskip
}
\fi
\newcommand{\ep}{\varepsilon}
\newcommand{\eps}{\varepsilon}
\makeatother
\title[Equivalence of intrinsic and extrinsic area bounds for minim]{Equivalence of intrinsic and extrinsic area bounds for minimal surfaces}
\author{Enric Florit-Simon}
\date{Source: arXiv:2605.06468v1 (CC BY 4.0)}
\begin{document}
\maketitle

\noindent\emph{Source and licence.} Equivalence of intrinsic and extrinsic area bounds for minimal surfaces. Version arXiv:2605.06468v1 (CC BY 4.0), distributed under the Creative Commons Attribution 4.0 International licence, as recorded in the arXiv licence metadata for this exact version and shown on the abstract page https://arxiv.org/abs/2605.06468v1. Full rights notice: licenses/arxiv-2605.06468-CC-BY-4.0.txt.\par\smallskip

\noindent\textit{Editorial scope.} Counted unit: the Lemma whose source label is \texttt{lem:chordarc} in the article (source0.tex lines 455--458, source label \texttt{lem:chordarc}), delivered together with the article's own complete original proof of it (source0.tex lines 459--479, one \texttt{proof} environment of 21 lines). No mathematics is written, repaired or re-derived: statement and proof are the article's own text line for line. Required context retained uncounted: rmk:trivbound (lines 388--390); lem:monformula (lines 425--431); eq:monfor (lines 427--430); rmk:lowdens (lines 451--454). Every definition, display and label of those lines is retained unchanged. Omitted from this package and not redistributed: 1--387, 391--424, 431--450, 480--526. Nothing inside a retained excerpt is omitted or altered: every retained body is delivered exactly as the article's lines are written, so no omission receipt of any kind accompanies this package. Citations: the retained excerpts carry no citation command at all, so no bibliography entry of the article is transcribed and no reference is redistributed. Reference adaptation: none was needed and none was made; every reference inside the delivered unit resolves inside it, so no unresolved reference or citation is produced. Printed slips kept literal: no printed word, symbol, hypothesis or number of the counted statement or of its proof was altered. Layout and class: the article's own class is \texttt{article} with packages including tocbibind, appendix, inputenc, graphicx, color, enumitem, mathtools, cancel, babel, hyperref, amsmath, amssymb, amsfonts, mathrsfs; the wrapper is the lane's pinned \texttt{amsart} editorial wrapper at 10pt on A4, which restates the theorem-like environments the retained text needs and adds the article's own macro definitions that the retained text needs (\texttt{\textbackslash R}, \texttt{\textbackslash ep}, \texttt{\textbackslash eps}), with extra wrapper packages (bm, bbm, dsfont, xspace, cancel, mathtools, cleveref). The delivered numbering is the unit's own. Assets: no retained span contains a figure environment, a graphics-inclusion command, a PDF-page inclusion, a table or a TikZ picture; no image file is copied into this package, the package declares no asset dependency and no image file is redistributed with it. Licence: version arXiv:2605.06468v1 of this article is distributed under the Creative Commons Attribution 4.0 International licence, as recorded in the arXiv licence metadata for this exact version and shown on the abstract page https://arxiv.org/abs/2605.06468v1. A full rights notice is delivered at licenses/arxiv-2605.06468-CC-BY-4.0.txt. Characters: every retained excerpt is pure ASCII, so the delivered engine, XeLaTeX with its default Latin Modern text font, reports no missing character.
\begin{remark}\label{rmk:trivbound}
    The inequality ${\bf M}_R^{\rm int}(\Sigma,p)\leq {\bf M}_R^{\rm ext}(\Sigma,p)$ is trivial and completely local, by the direct comparison $|p-q|\leq d_{\Sigma}(p,q)$ for any $p,q\in\Sigma\overset{i}{\hookrightarrow}\R^N$. Our proof consists in finding an asymptotic chord-arc estimate in the opposite direction, exploiting a monotonicity formula for the intrinsic area density.
\end{remark}

\begin{lemma}[Monotonicity formula]\label{lem:monformula}
    Let $\Sigma^d\overset{i}{\hookrightarrow}\R^N$ be an immersed submanifold with boundary $\partial\Sigma$ and $0\in\Sigma$. Let ${\rm H}_\Sigma$ denote its mean curvature vector. Then, given $0<R_1<R_2<{\rm dist}_\Sigma(0,\partial \Sigma)$, we have
    \begin{equation}\label{eq:monfor}
        {\bf M}_{R_2}^{\rm int}(\Sigma)-{\bf M}_{R_1}^{\rm int}(\Sigma)=
\int_{B_{R_2}^\Sigma\setminus B_{R_1}^\Sigma} \frac{(r_\Sigma(x)-|x|)+\frac{|x|}{2}(\frac{x}{|x|}-\nu_{B_{r_\Sigma(x)}^\Sigma}(x))^2}{r_\Sigma^{d+1}(x)}\,dx+\int_{R_1}^{R_2}\int_{B_{t}^\Sigma} \frac{x\cdot{\rm H}_\Sigma(x)}{t^{d+1}}\,dx\,dt\,.
\end{equation}
\end{lemma}

\begin{remark}\label{rmk:lowdens}
     In particular, if $\Sigma$ is minimal, recentering \Cref{lem:monformula} and sending $R_1\to 0$ we get the lower density bound
    $${\bf M}_R^{\rm int}(\Sigma,p)\geq |B_1^{\R^d}|=\omega_d\qquad \forall\, p\in \Sigma\quad\mbox{and}\quad R\in(0,{\rm dist}_\Sigma(0,\partial\Sigma)).$$
\end{remark}

\begin{lemma}[Density pinching implies chord-arc estimate]\label{lem:chordarc}
    Let $\Sigma^d\overset{i}{\hookrightarrow}\R^N$ be an immersed minimal submanifold with $0\in\Sigma$. Given $\eps\in(0,1)$, there exists $\delta=\delta(\eps,d)>0$ such that: If ${\bf M}_{8R}^{\rm int}(\Sigma)-{\bf M}_R^{\rm int}(\Sigma)\leq \delta$, then 
    $$|x|\leq r_\Sigma(x)\leq (1+\ep)|x|\qquad \mbox{in} \quad B_{4R}^\Sigma\setminus B_{2R}^\Sigma.$$
\end{lemma}

\begin{proof}
By scaling invariance we can assume that $R=1$.

    The inequality $|x|\leq r_\Sigma(x)$ is trivial (recall \Cref{rmk:trivbound}). Assume, for contradiction, that there were $x\in B_{4}^\Sigma\setminus B_{2}^\Sigma$ with $r_\Sigma(x)\geq (1+\eps)|x|$, thus with $r_\Sigma(x)-|x|\geq \frac{\eps}{1+\eps}r_\Sigma(x)\geq \frac{2\eps}{1+\eps}$. By the triangle inequality (both intrinsic and Euclidean), for every $y\in B_{\frac{\eps}{2(1+\eps)}}^\Sigma(x)\subset B_8^\Sigma\setminus B_1^\Sigma$ we still have
    \begin{equation}\label{eq:81tguqwp}
        r_\Sigma(y)-|y|\geq \big(r_\Sigma(x)-\frac{\eps}{2(1+\eps)}\big)-\big(|x|+\frac{\eps}{2(1+\eps)}\big)\geq \frac{\eps}{(1+\eps)}.
    \end{equation}
    On the other hand, \eqref{eq:monfor} shows that
    $$
    {\bf M}_8^{\rm int}(\Sigma)-{\bf M}_1^{\rm int}(\Sigma)\geq \int_{B_8^\Sigma\setminus B_1^\Sigma} \frac{r_\Sigma(y)-|y|}{r_\Sigma^{d+1}(y)}\,dy\geq \frac{1}{8^{d+1}}\int_{B_{\frac{\eps}{2(1+\eps)}}^\Sigma(x)} (r_\Sigma(y)-|y|)\,dy\,,
    $$
    thus combined with \eqref{eq:81tguqwp} we find
    $$
    {\bf M}_8^{\rm int}(\Sigma)-{\bf M}_1^{\rm int}(\Sigma) \geq \frac{\eps}{8^{d+1}(1+\eps)}|B_{\frac{\eps}{2(1+\eps)}}^\Sigma(x)|.
    $$
    By \Cref{rmk:lowdens} we can bound $|B_{\frac{\eps}{2(1+\eps)}}^\Sigma(x)|\geq \frac{\eps^d}{2^d(1+\eps)^d}|B_1^{\R^d}|$, thus
    $$
    {\bf M}_8^{\rm int}(\Sigma)-{\bf M}_1^{\rm int}(\Sigma) \geq \frac{\eps^{d+1}}{2^{4d+3}(1+\eps)^d}|B_1^{\R^d}|,
    $$
    reaching a contradiction for $\delta$ smaller than this value.
\end{proof}

\end{document}
